% !TEX encoding = UTF-8
% !TEX TS-program = pdflatex
% !TEX root = ../tesi.tex

%**************************************************************
\chapter{Introduzione}
\label{cap:introduzione}
%**************************************************************

Introduzione al contesto applicativo.\\

\noindent Esempio di utilizzo di un termine nel glossario \\
\gls{api}. \\

\noindent Esempio di citazione in linea \\
\cite{site:agile-manifesto}. \\

\noindent Esempio di citazione nel pie' di pagina \\
citazione\footcite{womak:lean-thinking} \\

%**************************************************************
\section{L'azienda AzzurroDigitale srl}

Lo stage è stato svolto nell'azienda AzzurroDigitale srl situata nella zona industriale di Padova. AzzurroDigitale nasce nel 2015 quando tre giovani padovani (Carlo Pasqualetto, Jacopo Pertile e Antonio Fornari) fondano la startup, puntando fortemente nelle nuove emergenti tecnologie che il mercato offriva. Come primo cliente, fu l'azienda Electrolux che grazie a una forte attività collaborazione, fu sviluppata una piattaforma per la gestione degli operai denominata AWMS, che tutt’ora continua a ricevere miglioramenti e a crescere. Dopo il successo ottenuto con la collaborazione con Electrolux, l'azienda capisce che il mercato delle aziende manifatturiere è la nicchia sulla quale puntare soprattutto grazie al momento storico della digital transformation.

\begin{figure}[htbp]
	\begin{center}
		\includegraphics[height=3cm]{Logo_azzurrodigite.png}
	\end{center}
\end{figure}

Oggi AzzurroDigitale offre servizi di industrial digital transformation, workforce management e people empowerment, con l'obbiettivo comune di aiutare le aziende manifatturiere a migliorare e implementare i loro processi grazie alle tecnologie, non intese come sostitutive all’uomo, ma bensì come mezzi che abilitano le persone a lavorare nel miglior modo possibile, massimizzando lo sforzo lavorativo.\\

%**************************************************************
\section{L'idea}

Introduzione all'idea dello stage.

%**************************************************************
\section{Organizzazione del testo}

\begin{description}
    \item[{\hyperref[cap:processi-metodologie]{Il secondo capitolo}}] descrive ...
    
    \item[{\hyperref[cap:descrizione-stage]{Il terzo capitolo}}] approfondisce ...
    
    \item[{\hyperref[cap:analisi-requisiti]{Il quarto capitolo}}] approfondisce ...
    
    \item[{\hyperref[cap:progettazione-codifica]{Il quinto capitolo}}] approfondisce ...
    
    \item[{\hyperref[cap:verifica-validazione]{Il sesto capitolo}}] approfondisce ...
    
    \item[{\hyperref[cap:conclusioni]{Nel settimo capitolo}}] descrive ...
\end{description}

Riguardo la stesura del testo, relativamente al documento sono state adottate le seguenti convenzioni tipografiche:
\begin{itemize}
	\item gli acronimi, le abbreviazioni e i termini ambigui o di uso non comune menzionati vengono definiti nel glossario, situato alla fine del presente documento;
	\item per la prima occorrenza dei termini riportati nel glossario viene utilizzata la seguente nomenclatura: \emph{parola}\glsfirstoccur;
	\item i termini in lingua straniera o facenti parti del gergo tecnico sono evidenziati con il carattere \emph{corsivo}.
\end{itemize}